\chapter{مدیریت پردازش‌ها (Processes)}

سیستم‌عامل‌های مدرن و در رأس آن‌ها لینوکس، ماهیتی چندوظیفه‌ای (\en{Multitasking}) دارند؛ بدین معنا که با ‌جابه‌جایی بسیار سریع میان برنامه‌های در حال اجرا، توهم اجرای هم‌زمان چندین عملیات را ایجاد می‌کنند. هسته لینوکس (\en{Kernel}) این نقش را از طریق مدیریت «پردازش‌ها» (\en{Processes}) به سرانجام می‌رساند. پردازش‌ها در واقع نمونه‌های در حال اجرایی از برنامه‌ها هستند که توسط هسته برای تخصیص بهینه منابع پردازشی سازماندهی می‌شوند.

در سناریوهایی که یک برنامه از پاسخ‌گویی باز می‌ماند یا نرخ بهره‌وری سیستم کاهش می‌یابد، درک عمیق از سازوکار پردازش‌ها و تسلط بر مدیریت آن‌ها اهمیتی حیاتی پیدا می‌کند. این فصل به معرفی ابزارهای توانمندی در محیط خط فرمان می‌پردازد که امکان پایش مستمر فعالیت برنامه‌ها و مداخله در فرآیند اجرای پردازش‌های مخرب یا بی‌پاسخ را فراهم می‌سازند.

در ادامه، فرامین کلیدی زیر را تحلیل خواهیم کرد:

\begin{itemize}
  \item \cmd{ps}: نمایش وضعیت لحظه‌ای (\en{Snapshot}) پردازش‌های فعال.
  \item \cmd{top}: نمایش پویا و بلادرنگ منابع سیستم و رتبه‌بندی پردازش‌ها.
  \item \cmd{jobs}: فهرست کردن وظایف (\en{Jobs}) در حال اجرا در نشست جاری پوسته.
  \item \cmd{bg}: انتقال یک پردازش متوقف‌شده به پس‌زمینه (\en{Background}).
  \item \cmd{fg}: انتقال یک پردازش از پس‌زمینه به پیش‌زمینه (\en{Foreground}).
  \item \cmd{kill}: ارسال سیگنال‌های کنترلی برای خاتمه دادن به یک پردازش.
  \item \cmd{killall}: خاتمه دادن به تمامی پردازش‌های هم‌نام.
  \item \cmd{nice}: تنظیم اولویت (\en{Priority}) یک برنامه پیش از اجرا.
  \item \cmd{renice}: اصلاح اولویت پردازش‌های در حال اجرا.
  \item \cmd{nohup}: اجرای پردازش به‌گونه‌ای که پس از بستن ترمینال نیز ادامه یابد.
  \item \cmd{halt/poweroff/reboot}: فرامین سیستمی جهت راه‌اندازی مجدد سیستم، خاموش‌سازی سیستم و توقف کامل سیستم.
  \item \cmd{shutdown}: خاموش‌سازی یا راه‌اندازی مجدد با قابلیت زمان‌بندی و اطلاع‌رسانی.
\end{itemize}

\section{سازوکار عملکرد پردازش‌ها}

هنگام بوت شدن سیستم، هسته لینوکس (\en{Kernel}) فعالیت‌های بنیادین خود را در قالب چندین پردازش آغاز کرده و متعاقباً برنامه \cmd{init} را اجرا می‌کند. این برنامه که نخستین پردازش سیستم و والد همهٔ پردازش‌های دیگر است، واحد \cmd{systemd} را فراخوانی می‌کند که مسئولیت راه اندازی و مدیریت تمام سرویس‌های سیستم را بر عهده دارد. بسیاری از این سرویس‌ها در قالب سرویس‌های پس‌زمینه یا همان دیمن (\en{Daemon}) پیاده‌سازی می‌شوند؛ برنامه‌هایی که در پس‌زمینه (\en{Background}) مستقر شده و بدون نیاز به تعامل مستقیم با کاربر، وظایف خود را انجام می‌دهند. از این رو، حتی پیش از احراز هویت کاربر، ده‌ها پردازش عملیاتی در حال اجرا هستند.

توانایی یک برنامه در فراخوانی و اجرای برنامه‌ای دیگر، منجر به شکل‌گیری یک ساختار سلسله‌مراتبی می‌شود؛ به گونه‌ای که یک فرآیند والد (\en{Parent Process})، یک یا چند فرآیند فرزند (\en{Child Process}) را ایجاد می‌کند. هسته لینوکس برای پایش و کنترل این ساختار شجره‌ای، شناسنامه دقیقی از اطلاعات هر فرآیند را در حافظه نگه‌داری می‌کند.

کلیدی‌ترین فاکتورهایی که هسته برای مدیریت هر پردازش ثبت می‌کند، عبارتند از:

\begin{itemize}
  \item \textbf{شناسه پردازش (\en{PID - Process ID}):} یک عدد منحصربه‌فرد که به هر پردازش الصاق می‌شود. این شناسه‌ها به صورت صعودی تخصیص یافته و همواره پردازش \cmd{init} (یا \cmd{systemd}) دارای \en{PID 1} است.
  \item \textbf{مدیریت حافظه:} هسته به‌طور دقیق میزان فضای آدرس‌دهی اختصاص‌یافته به هر پردازش را رصد می‌کند.
  \item \textbf{وضعیت پردازش (\en{Process State}):} وضعیت آمادگی یا تعلیق هر پردازش برای دسترسی به پردازنده توسط هسته پایش می‌شود.
  \item \textbf{اطلاعات مالکیت و امنیت:} مشابه فایل‌ها، هر پردازش دارای یک مالک (\en{Owner})، شناسه کاربری مؤثر (\en{Effective User ID}) و سایر ویژگی‌های امنیتی است که حدود دسترسی آن را تعیین می‌کند.
\end{itemize}

\section{دستور ps: مشاهده پردازش‌ها}

یکی از پرکاربردترین ابزارها جهت مشاهده و تحلیل پردازش‌های فعال در سیستم، فرمان \cmd{ps} (مخفف \en{Process Status}) است. این برنامه از انعطاف‌پذیری بالایی برخوردار است و گزینه‌های متعددی را می‌پذیرد؛ اما در ساده‌ترین حالت، بدون هیچ آرگومانی اجرا می‌شود:

\begin{latin}
\begin{lstlisting}
[me@linuxbox ~]$ ps
  PID TTY          TIME CMD
 5198 pts/1    00:00:00 bash
10129 pts/1    00:00:00 ps
\end{lstlisting}
\end{latin}

همان‌طور که در خروجی فوق مشهود است، فرمان \cmd{ps} در حالت پیش‌فرض تنها پردازش‌های مرتبط با نشست (\en{Session}) ترمینال جاری را فهرست می‌کند. در این مثال، تنها دو پردازش شناسایی شده‌اند: فرآیند \cmd{bash} (پوسته در حال اجرا) و خود فرآیند \cmd{ps} که برای استخراج این اطلاعات فراخوانی شده است.

\textbf{تحلیل ستون‌های خروجی در گزارش پردازش‌ها:}

\begin{itemize}
  \item \textbf{شناسه پردازش (\en{PID}):} یک شناسه عددی منحصربه‌فرد که توسط هسته به هر پردازش تخصیص می‌یابد. این شناسه مرجع اصلی برای مدیریت، پایش یا ارسال سیگنال‌های خاتمه به یک فرآیند محسوب می‌شود.
  \item \textbf{ترمینال کنترل‌کننده (\en{TTY}):} این ستون نشان‌دهنده دستگاه ترمینالی است که پردازش به آن وابسته است. عبارت \file{pts/1} (مخفف \en{Pseudo-Terminal Slave}) بیانگر آن است که پردازش در یک محیط شبیه‌ساز ترمینال (مانند محیط‌های گرافیکی یا نشست‌های راه دور) در حال اجراست. این اصطلاح بازمانده از دوران ترمینال‌های فیزیکی تله‌تایپ (\en{Teletype}) در سیستم‌عامل یونیکس است.
  \item \textbf{زمان مصرفی پردازنده (\en{TIME}):} مجموع مدت‌زمانی که پردازش مستقیماً از منابع پردازنده (\en{CPU}) استفاده کرده است. در خروجی فوق، مقدار \cmd{00:00:00} نشان می‌دهد که این پردازش‌ها منابع ناچیزی از پردازنده را به خود اختصاص داده‌اند.
  \item \textbf{فرمان اجرایی (\en{CMD}):} نام دقیق دستوری که منجر به ایجاد و فراخوانی پردازش مربوطه شده است.
\end{itemize}

با افزودن گزینه \cmd{x} به دستور \cmd{ps} (بدون استفاده از خط تیره یا \en{Dash} طبق قواعد نحوی \en{BSD})، می‌توان تمامی پردازش‌های متعلق به کاربر جاری را مشاهده کرد؛ این گزارش شامل پردازش‌هایی است که فاقد ترمینال کنترلی هستند.

\begin{latin}
\begin{lstlisting}
[me@linuxbox ~]$ ps x
  PID TTY      STAT   TIME COMMAND
 2799 ?        SsSl   0:00 /usr/libexec/bonobo-activation-server -ac
 2820 ?        Sl     0:01 /usr/libexec/evolution-data-server-1.10 --
15647 ?        Ss     0:00 /bin/sh /usr/bin/startkde
15751 ?        Ss     0:00 /usr/bin/ssh-agent /usr/bin/dbus-launch --
15754 ?        S      0:00 /usr/bin/dbus-launch --exit-with-session
15755 ?        Ss     0:01 /bin/dbus-daemon --fork --print-pid 4 -pr
15774 ?        Ss     0:02 /usr/bin/gpg-agent -s -daemon
15793 ?        S      0:00 start_kdeinit --new-startup +kcminit_start
15794 ?        Ss     0:00 kdeinit Running...
15797 ?        S      0:00 dcopserver -nosid
and many more...
\end{lstlisting}
\end{latin}

در این پیکربندی از خروجی، در ستون \en{TTY} علامت سوال (\cmd{?}) مشاهده می‌شود؛ این نماد نشان‌دهنده آن است که پردازش مورد نظر به هیچ ترمینال کنترل‌کننده‌ای متصل نیست. چنین وضعیتی معمولاً برای سرویس‌های پس‌زمینه (\en{Daemons}) یا فرآیندهای مربوط به رابط گرافیکی (\en{GUI}) رخ می‌دهد که مستقلاً اجرا می‌شوند.

به دلیل حجم بالای داده‌های تولیدی توسط \cmd{ps x} و طولانی بودن فهرست پردازش‌ها، توصیه می‌شود خروجی جهت پیمایش آسان‌تر به ابزار \cmd{less} پایپ (\en{Pipe}) شود:

\begin{latin}
\begin{lstlisting}
[me@linuxbox ~]$ ps x | less
\end{lstlisting}
\end{latin}

همچنین، بیشینه کردن (\en{Maximize}) ابعاد پنجره ترمینال به جلوگیری از شکستگی خطوط و بهبود خوانایی ستون‌های اطلاعاتی کمک شایانی می‌کند.

در خروجی دستور \cmd{ps} با گزینه‌های سبک \en{BSD} (مانند \cmd{ps aux} یا \cmd{ps x})، ستون \en{STAT} (مخفف \en{State} یا وضعیت) اطلاعات حیاتی درباره شرایط جاری هر فرآیند ارائه می‌دهد. این ستون با استفاده از کدهای حرفی، وضعیت عملیاتی پردازش را در چرخه حیات خود مشخص می‌کند که جزئیات آن‌ها در جدول زیر تدوین شده است.

\begin{table}[htbp]
\centering
\caption{وضعیت‌های پردازش (\en{Process States})}
\label{tab:process-states}
\begin{tabularx}{\textwidth}{c l X}
\toprule
\textbf{کد} & \textbf{مفهوم فنی} & \textbf{شرح کامل} \\
\midrule
\cmd{R} & در حال اجرا یا آمادهٔ اجرا (\en{Running or runnable}) & فرآیند به صورت فعال توسط پردازنده (\en{CPU}) در حال پردازش است یا در صف اجرا قرار دارد و منتظر تخصیص زمان‌بندی (\en{Time Slice}) بعدی است. \\
\addlinespace
\cmd{S} & خواب قابل‌وقفه (\en{Interruptible sleep}) & فرآیند موقتاً متوقف شده و منتظر وقوع یک رویداد خارجی (مانند ورودی کاربر یا دریافت داده از شبکه) است تا به وضعیت \cmd{R} بازگردد. \\
\addlinespace
\cmd{D} & خواب غیرقابل‌وقفه (\en{Uninterruptible sleep}) & فرآیند در انتظار تکمیل یک عملیات ورودی/خروجی (\en{I/O}) حیاتی — مانند دسترسی مستقیم به دیسک — است. در این حالت، فرآیند حتی با سیگنال‌های سیستمی نیز متوقف نمی‌شود تا یکپارچگی داده‌ها حفظ شود. \\
\addlinespace
\cmd{T} & متوقف‌شده (\en{Stopped}) & اجرای فرآیند توسط یک سیگنال کنترلی (مانند تعلیق توسط کاربر با \cmd{Ctrl+Z}) به حالت تعلیق درآمده است و تا دریافت سیگنال بازگشت، هیچ فعالیتی نخواهد داشت. \\
\addlinespace
\cmd{Z} & فرآیند زامبی (\en{Zombie / Defunct}) & فرآیند فرزندی که اجرای آن خاتمه یافته، اما رکورد آن همچنان در جدول فرآیندها باقی مانده است؛ زیرا فرآیند والد هنوز وضعیت خروج (\en{Exit Status}) آن را بازخوانی نکرده است. \\
\addlinespace
\cmd{<} & اولویت بالا (\en{high-priority}) & این فرآیند از نظر زمان‌بندی، اولویت بیشتری برای تصاحب منابع \en{CPU} دارد. در اصطلاح لینوکسی، این پردازش \en{``Less Nice''} تلقی می‌شود. \\
\addlinespace
\cmd{N} & اولویت پایین (\en{low-priority}) & فرآیندی که تنها در صورت آزاد بودن منابع به آن زمان پردازش تعلق می‌گیرد. این پردازش‌ها را \en{``Nice''} می‌نامند، زیرا منابع را برای سایر پردازش‌ها آزاد می‌گذارند. \\
\bottomrule
\end{tabularx}
\end{table}

علاوه بر کدهای پایه فوق، ممکن است کاراکترهای تکمیلی دیگری (نظیر \cmd{s} برای سرگروه نشست (\en{session leader}) یا \cmd{+} برای فرآیندهای پیش‌زمینه (\en{foreground})) نیز در کنار کد اصلی ظاهر شوند. برای بررسی دقیق‌تر این جزئیات فنی، مراجعه به مستندات رسمی سیستم از طریق فرمان \cmd{man ps} توصیه می‌شود.

یکی دیگر از ترکیب‌های بسیار رایج و کاربردی، استفاده از مجموعه گزینه‌های \cmd{aux} است. این الگو، خروجی دستور را گسترش داده و اطلاعات جامع‌تری از تمامی پردازش‌های سیستم ارائه می‌دهد:

\begin{latin}
\begin{lstlisting}
[me@linuxbox ~]$ ps aux
USER       PID %CPU %MEM    VSZ   RSS TTY      STAT START   TIME COMMAND
root         1  0.0  0.0   2136   644 ?        Ss   Mar05   0:31 init
root         2  0.0  0.0      0     0 ?        S<   Mar05   0:00 [kt]
root         3  0.0  0.0      0     0 ?        S<   Mar05   0:00 [mi]
root         4  0.0  0.0      0     0 ?        S<   Mar05   0:00 [ks]
root         5  0.0  0.0      0     0 ?        S<   Mar05   0:06 [wa]
root         6  0.0  0.0      0     0 ?        S<   Mar05   0:36 [ev]
root         7  0.0  0.0      0     0 ?        S<   Mar05   0:00 [kh]
and many more...
\end{lstlisting}
\end{latin}

به‌کارگیری گزینه‌های \cmd{aux} به هسته ابزار \cmd{ps} فرمان می‌دهد که پردازش‌های متعلق به تمامی کاربران سیستم را در خروجی درج کند. حذف خط تیره در این دستور باعث می‌شود \cmd{ps} مطابق با استانداردهای \en{BSD} رفتار کند. شایان ذکر است که نسخه لینوکسی این ابزار به‌گونه‌ای طراحی شده که قادر است شیوه‌های اجراییِ پیاده‌سازی‌های مختلف یونیکس را شبیه‌سازی نماید. پرکاربردترین گزینه‌های این سبک در جدول زیر تشریح شده‌اند:

\begin{table}[htbp]
\centering
\caption{گزینه‌های پرکاربرد \cmd{ps} به سبک \en{BSD}}
\label{tab:ps-bsd-options}
\begin{tabularx}{\textwidth}{c X}
\toprule
\textbf{گزینه} & \textbf{شرح عملکرد فنی} \\
\midrule
\cmd{x} & نمایش تمامی پردازش‌های متعلق به کاربر فعلی، از جمله پردازش‌های پس‌زمینه که فاقد ترمینال کنترل‌کننده (\en{TTY}) هستند. \\
\addlinespace
\cmd{ax} & نمایش تمامی فرآیندهای در حال اجرا در سطح سیستم، فارغ از اینکه متعلق به کدام کاربر باشند. \\
\addlinespace
\cmd{w} & نمایش کامل خط فرمان (\en{Command Line}) بدون قطع کردن یا برش انتهای عبارات طولانی؛ این گزینه برای واکاوی آرگومان‌های ورودی پردازش‌ها حیاتی است. \\
\addlinespace
\cmd{u} & تولید خروجی با جزئیات بالا (\en{User-oriented}) که ستون‌های جدیدی نظیر نام مالک، درصد اشغال پردازنده (\en{CPU\%}) و میزان مصرف حافظهٔ رم (\en{MEM\%}) را شامل می‌شود. \\
\bottomrule
\end{tabularx}
\end{table}

با ترکیب این گزینه‌ها در قالب \cmd{aux}، خروجی نهایی شامل ستون‌های آماری پیشرفته‌تری می‌گردد که تصویری شفاف از تخصیص منابع و وضعیت عملیاتی هر فرآیند در اختیار مدیر سیستم قرار می‌دهد.

\begin{table}[htbp]
\centering
\caption{تحلیل ستون‌های خروجی \cmd{ps} به سبک \en{BSD}}
\label{tab:ps-bsd-columns}
\begin{tabularx}{\textwidth}{l X}
\toprule
\textbf{ستون} & \textbf{مفهوم فنی و کاربردی} \\
\midrule
\en{USER} & نام کاربری یا شناسه کاربری (\en{UID}) مالک فرآیند. \\
\addlinespace
\en{\%CPU} & میزان بهره‌گیری از توان پردازشی هسته‌های \en{CPU} به‌صورت درصدی. \\
\addlinespace
\en{\%MEM} & درصد حافظهٔ فیزیکی: نسبت \en{RSS} فرآیند به کل حافظهٔ فیزیکی (\en{RAM}) سیستم، که به‌صورت درصدی نمایش داده می‌شود. این مقدار نشان‌دهندهٔ سهم یک فرآیند از کل حافظهٔ فیزیکی در دسترس است. \\
\addlinespace
\en{VSZ} & اندازه حافظه مجازی (\en{Virtual memory size}): کل فضای آدرس‌دهی‌شده‌ای که فرآیند می‌تواند به آن دسترسی داشته باشد. این مقدار نشان‌دهندهٔ مصرف واقعی حافظه نیست، بلکه اندازهٔ کل فضایی است که فرآیند برای خود رزرو کرده است. \\
\addlinespace
\en{RSS} & \en{Resident Set Size}: مقدار واقعی حافظهٔ فیزیکی (\en{RAM}) که در لحظه در اختیار فرآیند است. این مقدار شامل حافظهٔ مبادله‌شده (\en{Swap}) نمی‌شود و معیار کلیدی برای سنجش مصرف حافظهٔ فیزیکی توسط یک فرآیند محسوب می‌شود. \\
\addlinespace
\en{START} & لحظه آغاز به کار فرآیند؛ چنانچه از زمان شروع بیش از ۲۴ ساعت گذشته باشد، تاریخ جایگزین ساعت می‌شود. \\
\addlinespace
\en{TIME} & مجموع کل زمانی که فرآیند به‌طور مؤثر از پردازنده استفاده کرده است (\en{Cumulative CPU Time}). \\
\bottomrule
\end{tabularx}
\end{table}

برای تحلیل دقیق وضعیت یک فرآیند خاص، نیازی به مرور تمام فهرست سیستم نیست؛ شما می‌توانید با استفاده از شناسه پردازش (\en{PID}) به عنوان آرگومان، اطلاعات را محدود به همان مورد نمایید.

\textbf{مثال عملیاتی:}

\begin{latin}
\begin{lstlisting}
[me@linuxbox ~]$ ps uw 44719
USER   PID %CPU %MEM   VSZ  RSS TTY      STAT START   TIME COMMAND
me   44719  0.0  0.0 13480 6492 pts/1    S    15:57   0:00 bash
\end{lstlisting}
\end{latin}

در این نمونه، دستور \cmd{ps} با گزینه‌های \cmd{uw} (برای دریافت خروجی مفصل و عریض) برای \en{PID 44719} اجرا شده است. خروجی حاصل، جزئیات کاملی از پردازش \cmd{bash} شامل مالکیت، میزان مصرف دقیق منابع و وضعیت عملیاتی (\en{STAT}) را ارائه می‌دهد. این شیوه، ابزاری قدرتمند برای مدیران سیستم جهت عیب‌یابی (\en{Troubleshooting}) و پایش رفتارهای غیرعادی در پردازش‌های حساس محسوب می‌شود.

\section{دستور top: مشاهده پویای پردازش‌ها}

در حالی که فرمان \cmd{ps} تصویری ایستا (\en{Snapshot}) از وضعیت سیستم در لحظه‌ای خاص ارائه می‌دهد، ابزار \cmd{top} راهکاری قدرتمند برای نظارت مستمر، زنده و تعاملی بر فرآیندها فراهم می‌کند.

\begin{latin}
\begin{lstlisting}
[me@linuxbox ~]$ top
\end{lstlisting}
\end{latin}

با اجرای این دستور، محیطی به‌روز‌شونده پدیدار می‌شود که به‌صورت پیش‌فرض هر سه ثانیه یک‌بار نوسانات سیستم را پایش می‌کند. فرآیندها در این محیط بر اساس شاخص‌های مصرف منابع (عمدتاً میزان اشغال پردازنده) مرتب می‌شوند؛ به همین سبب آن را \en{top} می‌نامند، چرا که فرآیندهای پرمصرف همواره در صدر فهرست قرار می‌گیرند.

ساختار خروجی \cmd{top} به دو بخش کلیدی تقسیم می‌شود:

\begin{itemize}
  \item \textbf{ناحیه خلاصه وضعیت سیستم (\en{System Summary}):} در سطور بالایی، آمارهای کلان سیستم شامل زمان فعالیت (\en{Uptime})، تعداد کاربران متصل، میانگین بار سیستم (\en{Load Average})، وضعیت تخصیص حافظهٔ فیزیکی و فضای مبادله (\en{Swap}) نمایش داده می‌شود.
  \item \textbf{جدول فرآیندها (\en{Process Table}):} در بخش پایینی، جزییات عملیاتی هر فرآیند به‌صورت ستونی درج شده است. این جدول با هر بار تازه‌سازی (\en{Refresh})، تغییرات در مصرف منابع و وضعیت فرآیندها را به‌صورت آنی منعکس می‌کند.
\end{itemize}

در اینجا نمونه‌ای از خروجی این ابزار آورده شده است:

\begin{latin}
\begin{lstlisting}
top - 14:59:20 up 6:30,  2 users,  load average: 0.07, 0.02, 0.00
Tasks: 109 total,   1 running, 106 sleeping,   0 stopped,   2 zombie
Cpu(s):  0.7%us,  1.0%sy,  0.0%ni, 98.3%id,  0.0%wa,  0.0%hi,  0.0%si
Mem:   319496k total,  314860k used,    4636k free,   19392k buff
Swap:  875500k total,  149128k used,   726372k free,  114676k cach

  PID USER      PR  NI  VIRT  RES  SHR S %CPU %MEM    TIME+ COMMAND
 6244 me        39  19 31752 3124 2188 S  6.3  1.0  16:24.42 trackerd
11071 me        20   0  2304 1092  840 R  1.3  0.3   0:00.14 top
 6180 me        20   0  2700 1100  772 S  0.7  0.3   0:03.66 dbus-dae
 6321 me        20   0 20944 7248 6560 S  0.7  2.3   2:51.38 multiloa
 4955 root      20   0  104m 9668 5776 S  0.3  3.0   2:19.39 Xorg
    1 root      20   0  2976  528  476 S  0.0  0.2   0:03.14 init
    2 root      15  -5     0    0    0 S  0.0  0.0   0:00.00 kthreadd
    3 root      RT  -5     0    0    0 S  0.0  0.0   0:00.00 migratio
    4 root      15  -5     0    0    0 S  0.0  0.0   0:00.72 ksoftirq
    5 root      RT  -5     0    0    0 S  0.0  0.0   0:00.04 watchdog
    6 root      15  -5     0    0    0 S  0.0  0.0   0:00.42 events/0
    7 root      15  -5     0    0    0 S  0.0  0.0   0:00.06 khelper
   41 root      15  -5     0    0    0 S  0.0  0.0   0:01.08 kblockd/
   67 root      15  -5     0    0    0 S  0.0  0.0   0:00.00 kseriod
  114 root      20   0     0    0    0 S  0.0  0.0   0:01.62 pdflush
  116 root      15  -5     0    0    0 S  0.0  0.0   0:02.44 kswapd0
\end{lstlisting}
\end{latin}

این ابزار نه‌تنها برای مشاهده، بلکه برای مدیریت تعاملی نیز به کار می‌رود؛ برای مثال با فشردن کلید \cmd{k} می‌توانید یک فرآیند را متوقف کنید یا با کلید \cmd{M} فهرست را بر اساس مصرف حافظهٔ فیزیکی (\en{RAM}) مرتب نمایید.

بخش فوقانی خروجی فرمان \cmd{top} حاوی داده‌های حیاتی از وضعیت سیستم است که در ادامه به تشریح دقیق آن‌ها پرداخته شده است:

\begin{table}[htbp]
\centering
\caption{فیلدهای اطلاعاتی در بخش فوقانی دستور \cmd{top}}
\label{tab:top-fields}
\begin{tabularx}{\textwidth}{c l X}
\toprule
\textbf{ردیف} & \textbf{فیلد/شناسه} & \textbf{شرح مفاهیم فنی} \\
\midrule
\multirow{5}{*}{\cmd{1}} & \cmd{top} \cmd{14:59:20} & نام ابزار در حال اجرا به همراه زمان فعلی سیستم (ساعت دقیق محلی). \\
\addlinespace
 & \cmd{up 6:30} & مدت زمان فعالیت (\en{Uptime}): بیانگر فاصله زمانی سپری شده از آخرین بوت یا راه‌اندازی مجدد سیستم؛ در این سناریو سیستم ۶ ساعت و ۳۰ دقیقه در وضعیت عملیاتی بوده است. \\
\addlinespace
 & \cmd{2 users} & تعداد نشست‌های (\en{Sessions}) فعال کاربران که در حال حاضر به سیستم متصل هستند. \\
\addlinespace
 & \cmd{load average} & میانگین بار سیستم: این سه عدد نشان‌دهنده میانگین تعداد پردازش‌های در صف انتظار برای \en{CPU} در بازه‌های زمانی ۱، ۵ و ۱۵ دقیقه اخیر هستند. به‌طور کلی، مقادیر کمتر از تعداد هسته‌های پردازنده، نشان‌دهنده عملکرد روان سیستم است. \\
\midrule
\cmd{2} & \cmd{Tasks} & خلاصه وضعیت فرآیندها: پایش کمّی تمامی پردازش‌ها در حالات مختلف (\en{Running, Sleeping, Stopped, Zombie}). \\
\midrule
\multirow{6}{*}{\cmd{3}} & \cmd{Cpu(s)} & تفکیک مصرف پردازنده: وضعیت مصرف توان پردازشی بر حسب درصد برای انواع فعالیت‌ها: \\
\addlinespace
 & \cmd{0.7\%us} & سهم فرآیندهای کاربر (\en{User Space}). \\
\addlinespace
 & \cmd{1.0\%sy} & سهم فرآیندهای هسته (\en{Kernel Space}). \\
\addlinespace
 & \cmd{0.0\%ni} & سهم پردازش‌های با اولویت تغییر یافته (\en{Nice}). \\
\addlinespace
 & \cmd{98.3\%id} & میزان بیکاری سیستم (\en{Idle}). \\
\addlinespace
 & \cmd{0.0\%wa} & زمان انتظار برای عملیات ورودی/خروجی (\en{I/O Wait}). \\
\midrule
\cmd{4} & \cmd{Mem} & آمار حافظهٔ فیزیکی: گزارش دقیق میزان کل، مصرف‌شده، آزاد و حافظه تخصیص یافته به بافرها. \\
\midrule
\cmd{5} & \cmd{Swap} & آمار فضای مبادله‌ (\en{Swap}): وضعیت فضای مبادله (\en{Swap}) روی دیسک، شامل میزان کل، استفاده‌شده، فضای آزاد و میزان داده‌های کش‌شده. \\
\bottomrule
\end{tabularx}
\end{table}

فرمان \cmd{top} یک ابزار تعاملی و قدرتمند برای پایش لحظه‌ای فرآیندها در محیط خط فرمان است که از دستورهای کلیدی متعددی برای کنترل رابط کاربری بهره می‌برد. دو مورد از حیاتی‌ترین این فرامین عبارتند از:

\begin{itemize}
  \item \cmd{h}: فراخوانی صفحه راهنمای تعاملی جهت مشاهده تمامی کلیدهای کنترلی.
  \item \cmd{q}: خروج فوری از محیط برنامه و بازگشت به پوسته.
\end{itemize}

اگرچه محیط‌های دسکتاپ رایج لینوکس (مانند \en{GNOME} و \en{KDE}) ابزارهای گرافیکی مشابهی — نظیر \en{Task Manager} در ویندوز — را برای پایش سیستم ارائه می‌دهند، اما \cmd{top} دارای مزایای راهبردی قابل‌توجهی است. این ابزار به دلیل سرعت پاسخ‌گویی بالا و کمترین میزان مصرف منابع سیستم (\en{Low Overhead})، برای عیب‌یابی فرآیندها، به‌ویژه در شرایطی که سیستم با افت عملکرد و کندی مواجه است، گزینه‌ای بی‌بدیل محسوب می‌شود. از منظر مهندسی سیستم، یک ابزار نظارتی هرگز نباید خود به بخشی از بار اضافی و عامل کندی سیستم تبدیل شود.

\section{کنترل پردازش‌ها}

اکنون که با شیوه‌های پایش و نظارت بر فرآیندها آشنا شده‌اید، گام بعدی یادگیری فنون کنترل و هدایت آن‌هاست. برای اجرای این تمرین‌ها، از برنامه‌ای سبک به نام \cmd{xlogo} استفاده خواهیم کرد. این ابزار، یک برنامهٔ نمونه در سیستم پنجره‌بندی \en{X} (\en{X Window System}) — که زیرساخت گرافیکی سنتی در سیستم‌های شبه‌یونیکس محسوب می‌شود — است و صرفاً پنجره‌ای حاوی لوگوی \en{X} را نمایش می‌دهد.

برای شروع، این برنامه را از طریق ترمینال فراخوانی کنید:

\begin{latin}
\begin{lstlisting}
[me@linuxbox ~]$ xlogo
\end{lstlisting}
\end{latin}

پس از اجرای دستور، پنجره کوچکی مزین به لوگوی \en{X} ظاهر می‌شود. چنانچه این ابزار در توزیع شما نصب نیست، می‌توانید از ویرایشگرهای متن نظیر \cmd{gedit} یا \cmd{kwrite} به عنوان جایگزین استفاده کنید. با تغییر ابعاد پنجره و مشاهده بازترسیم صحیح لوگو، از فعال بودن فرآیند اطمینان حاصل کنید.

نکته حائز اهمیت این است که پس از اجرای \cmd{xlogo} در محیط ترمینال، اعلان خط فرمان (\en{Prompt}) ناپدید شده و کنترل محیط موقتاً از دست می‌رود. این پدیده به این دلیل است که پوسته (\en{Shell}) در وضعیت انتظار برای خاتمه یافتن برنامه قرار می‌گیرد؛ رفتاری استاندارد که در تمامی فرامین پیشین نیز صادق بوده است. تنها با بستن پنجره برنامه، فرآیند مذکور پایان یافته و کنترل محیط مجدداً به پوسته واگذار می‌شود.

\bookfigure{img-03.png}{برنامه \cmd{xlogo}}

\subsection{خاتمه دادن به پردازش}

برای قطع اجرای یک فرآیند در حال اجرا، می‌توان از ترکیب کلیدهای کنترلی در محیط ترمینال استفاده کرد. برای درک عملی این مکانیزم، از برنامه \cmd{xlogo} بهره می‌گیریم:

\begin{itemize}
  \item ابتدا فرمان \cmd{xlogo} را در ترمینال اجرا نمایید تا پنجره برنامه ظاهر شود.
  \item سپس روی محیط ترمینال تمرکز (\en{Focus}) کنید.
  \item حالا کلیدهای ترکیبی \cmd{Ctrl+c} را همزمان فشار دهید.
\end{itemize}

با فشردن این کلیدها، یک سیگنال وقفه (\en{Interrupt Signal}) که در لینوکس با نام \en{SIGINT} شناخته می‌شود، به فرآیندِ پیش‌زمینه ارسال می‌گردد. این سیگنال از برنامه درخواست می‌کند تا بر اساس پروتکل‌های استاندارد، عملیات خود را متوقف کرده و خاتمه یابد. همان‌طور که مشاهده خواهید کرد، پنجره \cmd{xlogo} بی‌درنگ بسته شده و اعلان خط فرمان (\en{Prompt}) مجدداً در دسترس قرار می‌گیرد.

\begin{latin}
\begin{lstlisting}
[me@linuxbox ~]$ xlogo
[me@linuxbox ~]$
\end{lstlisting}
\end{latin}

اگرچه این روش در اکثر برنامه‌های تحت خط فرمان و ابزارهای گرافیکی ساده کاملاً مؤثر است، اما باید توجه داشت که برخی برنامه‌های پیچیده‌تر ممکن است این سیگنال را نادیده بگیرند. در چنین مواردی، باید از روش‌های تهاجمی‌تر یا ابزارهای مدیریت فرآیند پیشرفته‌تری استفاده نمود.

\subsection{انتقال پردازش به پس‌زمینه}

چنانچه تمایل دارید برنامه‌ای را بدون مسدود سازی محیط خط فرمان اجرا کنید، می‌توانید آن را به پس‌زمینه (\en{Background}) منتقل نمایید. در این سازوکار، ترمینال دارای یک محیط پیش‌زمینه (\en{Foreground}) است که دستورهای تعاملی را پردازش می‌کند و یک محیط پس‌زمینه که فرآیندهای غیرتعاملی را به‌صورت پنهان مدیریت می‌نماید.

برای فراخوانی یک برنامه در حالت پس‌زمینه، کافی است نویسهٔ \cmd{\&} (آمپرسند) را به انتهای دستور خود اضافه کنید:

\begin{latin}
\begin{lstlisting}
[me@linuxbox ~]$ xlogo &
[1] 28236
[me@linuxbox ~]$
\end{lstlisting}
\end{latin}

پس از اجرای این فرمان، پنجره \cmd{xlogo} باز می‌شود، اما برخلاف حالت استاندارد، اعلان خط فرمان (\en{Prompt}) بی‌درنگ جهت دریافت دستورهای جدید در دسترس شما قرار می‌گیرد. همچنین، پوسته پیامی حاوی دو شناسه کلیدی را چاپ می‌کند:

\begin{itemize}
  \item \cmd{[1]}: این عدد بیانگر شماره وظیفه (\en{Job Number}) است که توسط پوسته به این فرآیند اختصاص یافته است.
  \item \cmd{28236}: این عدد نشان‌دهنده شناسه فرآیند (\en{PID}) منحصربه‌فردی است که سیستم‌عامل برای \cmd{xlogo} در نظر گرفته است.
\end{itemize}

با اجرای دستور \cmd{ps} می‌توانید نمایی از فرآیندهای فعال در نشست جاری ترمینال خود را مشاهده کنید:

\begin{latin}
\begin{lstlisting}
[me@linuxbox ~]$ ps
  PID TTY          TIME CMD
10603 pts/1    00:00:00 bash
28236 pts/1    00:00:00 xlogo
28239 pts/1    00:00:00 ps
\end{lstlisting}
\end{latin}

همان‌طور که در خروجی فوق مشهود است، علاوه بر فرآیندهای مربوط به مفسر فرمان (\cmd{bash}) و خودِ دستور \cmd{ps}، فرآیند \cmd{xlogo} نیز با شناسه اختصاصی خود (\en{PID}) در فهرست دیده می‌شود که وضعیت فعال و در حال اجرای آن را تایید می‌کند.

قابلیت «کنترل وظایف» (\en{Job Control}) در پوسته، مکانیزمی داخلی برای مدیریت و نظارت بر دستوراتی است که در پس‌زمینه (\en{Background}) اجرا می‌شوند. با بهره‌گیری از دستور \cmd{jobs} می‌توانید فهرستی از این وظایف فعال را مشاهده کنید:

\begin{latin}
\begin{lstlisting}
[me@linuxbox ~]$ jobs
[1]+  Running                 xlogo &
\end{lstlisting}
\end{latin}

خروجی فوق بیانگر آن است که وظیفه‌ای با شماره اختصاصی \cmd{[1]}، که معادل اجرای دستور \cmd{xlogo \&} است، در وضعیت عملیاتی قرار دارد.

سیستم‌عامل لینوکس به شما اجازه می‌دهد چندین دستور را به‌طور هم‌زمان در پس‌زمینه مستقر کنید. با درج نویسه \cmd{\&} در انتهای هر دستور، هر پردازش به عنوان یک وظیفه مستقل در پس‌زمینه آغاز به کار می‌کند:

\begin{latin}
\begin{lstlisting}
me@linuxbox:~$ xlogo & gedit &
[1] 47211
[2] 47212
\end{lstlisting}
\end{latin}

در این نمونه، دو وظیفه جدید با شماره‌های ترتیبی \cmd{[1]} و \cmd{[2]} و شناسه‌های پردازش (\en{PID}) منحصربه‌فرد، در لایه پس‌زمینه سیستم اجرا شده‌اند.

\subsection{انتقال پردازش به پیش‌زمینه}

فرآیندهایی که در پس‌زمینه (\en{Background}) اجرا می‌شوند، از ورودی‌های استاندارد صفحه‌کلید در محیط ترمینال، از جمله سیگنال وقفه \cmd{Ctrl+C} مصون هستند. جهت بازگرداندن یک فرآیند به وضعیت پیش‌زمینه (\en{Foreground}) و برقراری تعامل مجدد با آن، از دستور \cmd{fg} استفاده می‌شود.

\textbf{نمونه عملیاتی:} ابتدا با استفاده از دستور \cmd{jobs} فهرست وظایف مستقر در پس‌زمینه را فراخوانی می‌کنیم:

\begin{latin}
\begin{lstlisting}
[me@linuxbox ~]$ jobs
[1]+  Running                 xlogo &
\end{lstlisting}
\end{latin}

خروجی نشان می‌دهد که پردازش \cmd{xlogo} با شماره وظیفه \cmd{[1]} در حال اجراست. برای انتقال آن به پیش‌زمینه، دستور \cmd{fg} را به‌همراه «شماره وظیفه» (\en{Jobspec}) اجرا می‌کنیم:

\begin{latin}
\begin{lstlisting}
[me@linuxbox ~]$ fg %1
xlogo
\end{lstlisting}
\end{latin}

با اجرای این فرمان، پردازش \cmd{xlogo} به پیش‌زمینه منتقل شده و کنترل ترمینال را در اختیار می‌گیرد؛ در این حالت می‌توانید با فشردن ترکیبی \cmd{Ctrl+C} به فعالیت آن خاتمه دهید. شایان ذکر است در صورتی که تنها یک وظیفه در پس‌زمینه وجود داشته باشد، استفاده از دستور \cmd{fg} به تنهایی و بدون قید کردن شماره وظیفه کفایت می‌کند.

\subsection{تعلیق پردازش}

در مواردی نیاز است که یک فرآیند را بدون خاتمه دادن به چرخه حیات آن، به‌طور موقت متوقف کنیم. این اقدام معمولاً گامی واسط برای انتقال یک پردازش از پیش‌زمینه (\en{Foreground}) به پس‌زمینه (\en{Background}) محسوب می‌شود. جهت دستیابی به این هدف، از ترکیب کلیدهای کنترلی \cmd{Ctrl+Z} استفاده می‌شود.

برای بررسی این سازوکار، مراحل زیر را دنبال کنید:

\begin{itemize}
  \item در خط فرمان، دستور \cmd{xlogo} را اجرا کنید.
  \item بلافاصله پس از اجرا، کلیدهای \cmd{Ctrl+Z} را بفشارید.
\end{itemize}

\begin{latin}
\begin{lstlisting}
[me@linuxbox ~]$ xlogo
[1]+  Stopped                 xlogo
[me@linuxbox ~]$
\end{lstlisting}
\end{latin}

با فشردن این کلیدها، فرآیند \cmd{xlogo} به وضعیت «متوقف‌شده» (\en{Stopped}) تغییر حالت می‌دهد. در این وضعیت، اجرای پردازش به‌صورت موقت به تعلیق درآمده، اما از حافظه سیستم خارج نمی‌شود. در این مرحله، پوسته با نمایش پیامی مبنی بر توقف وظیفه شماره \cmd{[1]}، کنترل خط فرمان (\en{Prompt}) را مجدداً به شما باز می‌گرداند. اکنون می‌توانید با بهره‌گیری از دستور \cmd{bg} (مخفف \en{Background})، اجرای این فرآیندِ معلق را در پس‌زمینه از سر بگیرید.

پس از تعلیق فرآیند \cmd{xlogo} با فشردن کلیدهای ترکیبی، می‌توانید با تلاش برای تغییر ابعاد پنجره آن، وضعیت توقف برنامه را احراز کنید؛ در این حالت پنجره برنامه کاملاً فاقد پاسخ‌دهی (\en{Non-responsive}) خواهد بود. برای بازگرداندن فرآیند به چرخه حیات، می‌توانید از دستور \cmd{fg} جهت اجرا در پیش‌زمینه و یا دستور \cmd{bg} جهت تداوم اجرا در پس‌زمینه استفاده کنید:

\begin{latin}
\begin{lstlisting}
[me@linuxbox ~]$ bg %1
[1]+ xlogo &
[me@linuxbox ~]$
\end{lstlisting}
\end{latin}

مشابه با دستور \cmd{fg}، در اینجا نیز در صورت وجود تنها یک وظیفه در فهرست، قید کردن «شماره وظیفه» (\en{Jobspec}) اختیاری است.

انتقال یک پردازش از پیش‌زمینه به پس‌زمینه، به‌ویژه زمانی که یک برنامه گرافیکی را از طریق ترمینال فراخوانی کرده‌اید اما نویسه \cmd{\&} را برای ارسال آن به پس‌زمینه فراموش نموده‌اید، بسیار کارآمد است. این تکنیک به شما اجازه می‌دهد بدون نیاز به بستن برنامه و از دست دادن نشست جاری، کنترل ترمینال را برای ادامه فعالیت‌های خود بازپس بگیرید.

اجرای نرم‌افزارهای مبتنی بر رابط کاربری گرافیکی (\en{GUI}) از طریق ترمینال، به دلایل فنی متعددی صورت می‌گیرد که آن را به مهارتی ارزشمند برای کاربران حرفه‌ای تبدیل می‌کند:

\begin{itemize}
  \item \textbf{دسترسی به برنامه‌های ثبت‌نشده در منوها:} ممکن است برخی برنامه‌ها در منوی برنامه‌های محیط دسکتاپ یا مدیر پنجره فهرست نشده باشند. در چنین مواردی، خط فرمان به‌عنوان مجرای مستقیم فراخوانی آن‌ها عمل می‌کند؛ همان‌طور که در مثال ابزار \cmd{xlogo} مشاهده کردیم.
  \item \textbf{عیب‌یابی و مشاهدهٔ پیام‌های خطا:} هنگام اجرای یک برنامهٔ گرافیکی از طریق ترمینال، تمام پیام‌های خطا، هشدارها و گزارش‌های خروجی آن در همان پنجرهٔ ترمینال نمایش داده می‌شوند. این ویژگی، فرآیند عیب‌یابی مشکلات مربوط به اجرا یا عملکرد برنامه را به‌شدت تسهیل می‌کند، زیرا توسعه‌دهندگان اغلب اطلاعات اشکال‌زدایی را از این طریق در اختیار کاربر قرار می‌دهند.
  \item \textbf{استفاده از گزینه‌ها و آرگومان‌های پیشرفته:} بسیاری از برنامه‌های گرافیکی، گزینه‌ها و پارامترهای خط فرمانی را می‌پذیرند که از طریق رابط گرافیکی قابل دسترسی نیستند. این گزینه‌ها می‌توانند برای اجرای برنامه با تنظیمات خاص، فعال‌سازی حالت‌های پیشرفته یا تغییر رفتار پیش‌فرض آن به کار روند.
\end{itemize}

\section{دستورات nice و renice: تغییر اولویت پردازش‌ها}

دستور \cmd{nice} جهت تخصیص یک «شاخص ملایمت» (\en{Niceness Value}) مشخص به فرآیند در هنگام شروع به‌کار آن استفاده می‌شود. این مقدار عددی در بازه‌ای میان \cmd{-20} (بالاترین اولویت دسترسی به \en{CPU}) تا \cmd{19} (کمترین اولویت) تعریف شده است؛ به این صورت که مقدار پیش‌فرض به‌صورت خودکار روی \cmd{0} تنظیم می‌شود.

به‌عنوان مثال، تصور کنید برنامه‌ای تحت عنوان \cmd{cpu-hog} در اختیار دارید که منابع پردازشی زیادی مصرف می‌کند و تمایل دارید آن را با اولویتی کمتر از سطح استاندارد اجرا کنید. در این صورت، می‌توانید دستور را به شیوه زیر وارد نمایید:

\begin{latin}
\begin{lstlisting}
[me@linuxbox ~]$ nice -n 10 cpu-hog
\end{lstlisting}
\end{latin}

در این فرمان، آرگومان \cmd{-n 10} به دستور \cmd{nice} ابلاغ می‌کند که فرآیند \cmd{cpu-hog} را با شاخص ملایمت \cmd{10} اجرا کند. این اقدام موجب می‌شود اولویت زمان‌بندی (\en{Scheduling}) این برنامه نسبت به حالت پیش‌فرض کاهش یافته و سهم کمتری از چرخه‌های پردازنده را به خود اختصاص دهد. به همین ترتیب، چنانچه برنامه‌ای حیاتی نظیر \cmd{must-run-fast} در اختیار داشته باشید که به اولویت بالاتری برای دسترسی به \en{CPU} نیاز دارد، می‌توانید در نقش «ابرکاربر» (\en{superuser}) آن را با استفاده از \cmd{sudo nice} و یک مقدار منفی برای شاخص ملایمت اجرا کنید:

\begin{latin}
\begin{lstlisting}
[me@linuxbox ~]$ sudo nice -n -10 must-run-fast
\end{lstlisting}
\end{latin}

شایان ذکر است که کاهش شاخص ملایمت و افزایش اولویت یک پردازش به‌ندرت ضرورت می‌یابد؛ چرا که این اقدام ممکن است فرآیندهای حیاتی سیستم را از منابع مورد نیاز محروم ساخته و منجر به ناپایداری شود. از این رو، اعمال چنین تغییراتی مستلزم دقت فنی و احتیاط کامل است.

دستور \cmd{renice} امکان تغییر اولویت زمان‌بندی فرآیندهایی را که پیش‌تر در سیستم آغاز به کار کرده‌اند، فراهم می‌سازد. با بهره‌گیری از \cmd{renice} قادر خواهید بود «شاخص ملایمت» (\en{Niceness}) یک پردازش در حال اجرا را بدون نیاز به توقف یا بازراه‌اندازی آن تغییر دهید. این فرمان نیز از همان بازه استاندارد \cmd{-20} تا \cmd{19} برای تعیین سطح ملایمت استفاده می‌کند. برای نمونه، اگر برنامه \cmd{cpu-hog} را اجرا کرده‌اید و قصد دارید شاخص ملایمت (\en{Niceness}) آن را افزایش دهید تا سهم کمتری از منابع را اشغال کند، می‌توانید مراحل زیر را دنبال کنید:

ابتدا با اجرای دستور \cmd{ps} در ترمینال، «شناسه فرآیند» (\en{PID}) مربوط به برنامه \cmd{cpu-hog} را شناسایی کنید:

\begin{latin}
\begin{lstlisting}
[me@linuxbox ~]$ ps
   PID TTY          TIME CMD
379087 pts/9    00:00:00 bash
379215 pts/9    00:00:00 cpu-hog
379223 pts/9    00:00:00 ps
\end{lstlisting}
\end{latin}

سپس با استفاده از دستور \cmd{renice} و قید کردن تراز جدید ملایمت به‌همراه \en{PID} مربوطه، تغییرات را اعمال نمایید:

\begin{latin}
\begin{lstlisting}
[me@linuxbox ~]$ renice -n 19 379215
\end{lstlisting}
\end{latin}

در این مثال، شاخص ملایمت بر روی عدد \cmd{19} (کمترین اولویت ممکن) تنظیم شده است. این پیکربندی تضمین می‌کند که فرآیند \cmd{cpu-hog} تنها زمانی مجاز به استفاده از چرخه‌های پردازنده باشد که هیچ فرآیند دیگری در صف انتظارِ سیستم وجود ندارد. این تکنیک رویکردی کارآمد برای مدیریت پویای منابع سیستم، بدون نیاز به متوقف کردن فرآیندهای جاری است.

\section{دستور kill: ارسال سیگنال به پردازش‌ها}

دستور \cmd{kill} ابزاری کلیدی در مدیریت فرآیندهاست که از طریق ارسال سیگنال، با پردازش‌ها ارتباط برقرار می‌کند. برخلاف تصور رایج، این دستور لزوماً باعث توقف آنی فرآیند نمی‌شود؛ بلکه پیامی (سیگنال) را مخابره می‌کند تا پردازش بر اساس نوع آن پیام، واکنش مشخصی نشان دهد. ساختار نحوی این دستور به شرح زیر است:

\begin{latin}
\begin{lstlisting}
kill [-signal] PID...
\end{lstlisting}
\end{latin}

چنانچه نوع سیگنال در دستور قید نشود، به‌صورت پیش‌فرض سیگنال \en{TERM} ارسال می‌گردد که از فرآیند می‌خواهد به‌طور منظم و ایمن به فعالیت خود پایان دهد.

برای استفاده از دستور \cmd{kill} کافی است «شناسه فرآیند» (\en{PID}) یا «شماره وظیفه» (\en{Jobspec}) را به عنوان آرگومان به دستور پاس دهید.

\textbf{مثال عملی:}

\begin{latin}
\begin{lstlisting}
[me@linuxbox ~]$ xlogo &
[1] 28401
[me@linuxbox ~]$ kill 28401
[1]+  Terminated              xlogo
\end{lstlisting}
\end{latin}

در این نمونه، ابتدا برنامهٔ \cmd{xlogo} در پس‌زمینه اجرا شد و سیستم شماره وظیفهٔ \cmd{[1]} و شمارهٔ \en{PID} برابر با \cmd{28401} را به آن اختصاص داد. سپس با اجرای دستور \cmd{kill 28401}، یک سیگنال خاتمهٔ استاندارد به آن ارسال کردیم که منجر به بسته شدن اصولی برنامه شد. توجه داشته باشید که می‌توانستیم به‌جای \en{PID}، از شناسه وظیفه (\en{Jobspec}) به صورت \cmd{\%1} نیز برای هدف قرار دادن این فرآیند استفاده کنیم.

سیگنال‌ها از بنیادی‌ترین روش‌های ارتباطی میان هسته سیستم‌عامل و برنامه‌ها محسوب می‌شوند. پیش از این، تأثیر سیگنال‌ها را در قالب میان‌برهای صفحه‌کلید تجربه کرده‌ایم:

\begin{itemize}
  \item \cmd{Ctrl+C}: سیگنال \en{SIGINT} (مخفف \en{Interrupt}) را ارسال می‌کند.
  \item \cmd{Ctrl+Z}: سیگنال \en{SIGTSTP} (مخفف \en{Terminal Stop}) را ارسال می‌کند.
\end{itemize}

برنامه‌ها به‌طور مداوم برای دریافت این سیگنال‌ها گوش‌به‌زنگ هستند و می‌توانند پس از دریافت، واکنش‌های برنامه‌ریزی‌شده‌ای نشان دهند. این قابلیت به نرم‌افزار اجازه می‌دهد پیش از خروج کامل از حافظه، عملیات‌های حیاتی نظیر ذخیره‌سازی داده‌های موقت، بستن فایل‌های باز یا آزادسازی منابع سیستمی را به انجام برساند. چنین رویکردی، مدیریت فرآیندها را از یک اقدام قهری به یک تعامل هماهنگ و منظم میان نرم‌افزار و سیستم تبدیل می‌کند.

\begin{table}[htbp]
\centering
\caption{سیگنال‌های متداول}
\label{tab:common-signals}
\begin{tabularx}{\textwidth}{c l X}
\toprule
\textbf{شماره} & \textbf{نام سیگنال} & \textbf{مفهوم و کاربرد فنی} \\
\midrule
\cmd{1} & \cmd{HUP} (\en{Hangup}) & سیگنال قطع ارتباط: این سیگنال که ریشه در دوران مودم‌های دیال‌آپ دارد، بیانگر قطع ارتباط ترمینالِ کنترل‌کننده است. هنگام بستن یک نشست (\en{Session}) ترمینال، این سیگنال به فرآیندهای پیش‌زمینه ارسال شده و موجب خاتمه آن‌ها می‌گردد. امروزه بسیاری از سرویس‌های سیستمی (\en{Daemon}) از این سیگنال برای بازخوانی فایل‌های پیکربندی (\en{Reload}) بدون نیاز به توقف کامل استفاده می‌کنند. \\
\addlinespace
\cmd{2} & \cmd{INT} (\en{Interrupt}) & سیگنال وقفه: عملکردی مشابه فشردن کلیدهای \cmd{Ctrl+C} دارد و به‌صورت مسالمت‌آمیز از برنامه درخواست می‌کند تا به فعالیت خود پایان دهد. \\
\addlinespace
\cmd{9} & \cmd{KILL} (\en{Kill}) & سیگنال نابودی: این سیگنال ماهیتی منحصربه‌فرد دارد؛ چرا که به جای ارسال به خودِ برنامه، مستقیماً توسط هسته (\en{Kernel}) دریافت شده و فرآیند را به‌سرعت نابود می‌کند. از آنجا که پردازش فرصتی برای ذخیره‌سازی داده‌ها یا پاک‌سازی منابع ندارد، این سیگنال باید صرفاً به عنوان آخرین راهکار استفاده شود. \\
\addlinespace
\cmd{15} & \cmd{TERM} (\en{Terminate}) & سیگنال خاتمه: سیگنال استاندارد و پیش‌فرض دستور \cmd{kill} است که از فرآیند می‌خواهد تا به‌صورت منظم و ایمن خاتمه یابد. \\
\addlinespace
\cmd{18} & \cmd{CONT} (\en{Continue}) & سیگنال ادامه: جهت از سرگیری اجرای فرآیندهایی که پیش‌تر با سیگنال‌های \en{STOP} یا \en{TSTP} معلق شده‌اند، به کار می‌رود. این سیگنال توسط دستورات \cmd{bg} و \cmd{fg} ارسال می‌شود. \\
\addlinespace
\cmd{19} & \cmd{STOP} (\en{Stop}) & سیگنال توقف اجباری: فرآیند را بدون خاتمه دادن، به‌طور کامل متوقف (\en{Suspend}) می‌کند. مشابه سیگنال \en{KILL}، این مورد نیز توسط هسته مدیریت شده و توسط برنامه قابل چشم‌پوشی یا مسدودسازی نیست. \\
\addlinespace
\cmd{20} & \cmd{TSTP} (\en{Terminal Stop}) & سیگنال توقف کاربر: سیگنالی است که هنگام فشردن کلیدهای \cmd{Ctrl+Z} صادر می‌شود. این سیگنال توسط برنامه دریافت می‌گردد و نرم‌افزار می‌تواند بر اساس منطق داخلی خود، آن را مدیریت کرده یا حتی نادیده بگیرد. \\
\bottomrule
\end{tabularx}
\end{table}

بیایید عملکرد دستور \cmd{kill} را بررسی کنیم:

\begin{latin}
\begin{lstlisting}
[me@linuxbox ~]$ xlogo &
[1] 13546
[me@linuxbox ~]$ kill -1 13546
[1]+  Hangup                xlogo
\end{lstlisting}
\end{latin}

در این نمونه، ابتدا برنامه \cmd{xlogo} را در پس‌زمینه فراخوانی کرده و سپس سیگنال \en{HUP} (شماره \cmd{1}) را به آن ارسال کردیم. فرآیند \cmd{xlogo} خاتمه یافته و مفسر فرمان (\en{Shell}) با پیامی تایید می‌کند که پردازش پس‌زمینه سیگنال \en{hangup} را دریافت کرده است (ممکن است برای مشاهده این پیام، نیاز باشد کلید \en{Enter} را مجدداً فشار دهید). شایان ذکر است که در ساختار این دستور، سیگنال‌ها را می‌توان هم با شماره و هم با نام (با یا بدون پیشوند \en{SIG}) ارجاع داد:

\begin{latin}
\begin{lstlisting}
[me@linuxbox ~]$ xlogo &
[1] 13601
[me@linuxbox ~]$ kill -INT 13601
[1]+  Interrupt             xlogo
[me@linuxbox ~]$ xlogo &
[1] 13608
[me@linuxbox ~]$ kill -SIGINT 13608
[1]+  Interrupt             xlogo
\end{lstlisting}
\end{latin}

شما می‌توانید با تکرار مثال‌های فوق، سایر سیگنال‌ها را نیز آزمایش کنید؛ همچنین به یاد داشته باشید که استفاده از شماره وظیفه (\en{Jobspec}) به جای \en{PID} در تمامی این موارد مجاز است.

در نهایت، توجه به امنیت سیستم ضروری است: فرآیندها همانند فایل‌ها دارای مالک (\en{Owner}) هستند. برای ارسال سیگنال به یک پردازش، شما باید مالک آن باشید یا با دسترسی «ابرکاربر» (\en{Superuser}) اقدام کنید.

افزون بر سیگنال‌های پیشین که عمدتاً در تعامل با دستور \cmd{kill} کاربرد دارند، مجموعه‌ای از سیگنال‌های دیگر نیز وجود دارند که غالباً توسط خودِ سیستم مدیریت می‌شوند. این موارد در جدول زیر تشریح شده‌اند:

\begin{table}[htbp]
\centering
\caption{سایر سیگنال‌های متداول}
\label{tab:other-signals}
\begin{tabularx}{\textwidth}{c l X}
\toprule
\textbf{شماره} & \textbf{نام سیگنال} & \textbf{مفهوم و کاربرد فنی} \\
\midrule
\cmd{3} & \cmd{QUIT} & این سیگنال که معمولاً با فشردن کلیدهای ترکیبی \cmd{Ctrl+\textbackslash} صادر می‌شود، از برنامه می‌خواهد تا فوراً متوقف شده و یک فایل \en{Core Dump} (گزارشی از وضعیت حافظه در لحظه توقف) جهت تحلیل و عیب‌یابی تولید کند. \\
\addlinespace
\cmd{11} & \cmd{SEGV} & خطای قطعه‌بندی یا نقض دسترسی (\en{Segmentation Violation}). این سیگنال زمانی از سوی هسته به پردازش ارسال می‌شود که برنامه قصد دسترسی غیرمجاز به بخشی از حافظه را داشته باشد؛ برای مثال، تلاش برای بازنویسی در ناحیه‌ای که صرفاً «فقط-خواندنی» (\en{Read-only}) تعریف شده است. \\
\addlinespace
\cmd{28} & \cmd{WINCH} & تغییر ابعاد پنجره (\en{Window Change}). این سیگنال زمانی صادر می‌شود که کاربر ابعاد پنجره ترمینال را تغییر دهد. برنامه‌های تعاملی و مبتنی بر متن نظیر \cmd{top}، \cmd{less} یا ویرایشگر \cmd{vim} با دریافت این سیگنال، رابط کاربری خود را متناسب با ابعاد جدید بازطراحی می‌کنند. \\
\bottomrule
\end{tabularx}
\end{table}

برای استخراج لیست جامعی از تمام سیگنال‌های مورد پشتیبانی در سیستم خود، می‌توانید از دستور زیر استفاده کنید:

\begin{latin}
\begin{lstlisting}
[me@linuxbox ~]$ kill -l
\end{lstlisting}
\end{latin}

این فرمان، تمامی سیگنال‌های موجود را به‌همراه شماره و نام اختصاصی آن‌ها نمایش می‌دهد. دسترسی به این فهرست برای عیب‌یابی دقیق و مدیریت پیشرفته فرآیندها در محیط‌های عملیاتی بسیار حائز اهمیت است.

\section{دستور nohup: مقاوم‌سازی پردازش در برابر سیگنال Hangup}

همان‌طور که پیش‌تر اشاره شد، بسیاری از ابزارهای خط فرمان به محض دریافت سیگنال \en{HUP} (که در نتیجه بسته شدن نشست یا قطع ارتباط ترمینال صادر می‌شود) خاتمه می‌یابند. جهت جلوگیری از این رفتار و تضمین تداوم اجرای برنامه، می‌توان از دستور \cmd{nohup} (مخفف \en{No Hangup}) استفاده کرد. این دستور، فرآیند را در برابر سیگنال \en{HUP} ایمن کرده و به آن اجازه می‌دهد حتی پس از خروج کاربر یا بسته شدن ترمینال، به فعالیت خود ادامه دهد.

برای درک بهتر، اجرای عادی \cmd{xlogo} را در نظر بگیرید:

\begin{latin}
\begin{lstlisting}
[me@linuxbox ~]$ xlogo
\end{lstlisting}
\end{latin}

در این حالت، با بستن پنجره ترمینال، فرآیند \cmd{xlogo} نیز بلافاصله متوقف می‌شود. اما چنانچه برنامه را با پیشوند \cmd{nohup} فراخوانی کنید:

\begin{latin}
\begin{lstlisting}
[me@linuxbox ~]$ nohup xlogo
\end{lstlisting}
\end{latin}

در این وضعیت، حتی با بستن پنجره ترمینال، \cmd{xlogo} در پس‌زمینه به اجرای خود ادامه خواهد داد. این رویکرد برای اجرای اسکریپت‌های محاسباتی سنگین یا برنامه‌های زمان‌بری که نباید با قطع ارتباط کاربر دچار اختلال شوند، یک راهکار استاندارد و کاربردی است. به‌طور پیش‌فرض، \cmd{nohup} خروجی‌های برنامه را در فایلی به نام \file{nohup.out} ذخیره می‌کند تا پس از بسته شدن ترمینال، امکان بررسی نتایج فراهم باشد.

\section{دستور killall: ارسال سیگنال به چندین پردازش}

جهت مخابره سیگنال به چندین فرآیند به‌طور هم‌زمان، دستور \cmd{killall} ابزاری بسیار کارآمد است. این فرمان، پردازش‌ها را بر اساس نام برنامه یا نام کاربری هدف قرار می‌دهد. ساختار نحوی (\en{Syntax}) این دستور به شرح زیر است:

\begin{latin}
\begin{lstlisting}
killall [-u user] [-signal] name...
\end{lstlisting}
\end{latin}

برای درک بهتر عملکرد این ابزار، چند نمونه از برنامه \cmd{xlogo} را فراخوانی کرده و سپس آن‌ها را به‌صورت گروهی متوقف می‌کنیم:

\begin{latin}
\begin{lstlisting}
[me@linuxbox ~]$ xlogo &
[1] 18801
[me@linuxbox ~]$ xlogo &
[2] 18802
[me@linuxbox ~]$ killall xlogo
[1]-  Terminated              xlogo
[2]+  Terminated              xlogo
\end{lstlisting}
\end{latin}

در این مثال، دو نسخه از \cmd{xlogo} در پس‌زمینه اجرا شدند. سپس با اجرای فرمان \cmd{killall xlogo}، سیگنال پیش‌فرض \en{TERM} به تمامی فرآیندهای هم‌نام ارسال شد و هر دو به‌طور هم‌زمان خاتمه یافتند.

مشابه با دستور \cmd{kill}، برای ارسال سیگنال به فرآیندهایی که در مالکیت حساب کاربری شما نیستند، برخورداری از امتیازات «ابرکاربر» (\en{Superuser}) الزامی است.

\section{خاموش‌سازی و بازراه‌اندازی سیستم}

فرایند خاموش کردن سیستم در لینوکس، شامل مجموعه‌ای از اقدامات هماهنگ است که با هدف خاتمه منظم تمامی پردازش‌ها و انجام عملیات حیاتی (نظیر همگام‌سازی یا \en{Sync} کردن سیستم‌های فایل) پیش از قطع جریان برق صورت می‌گیرد. برای مدیریت این وضعیت، چهار دستور اصلی در لینوکس تعبیه شده است: \cmd{halt}، \cmd{poweroff}، \cmd{reboot} و \cmd{shutdown}. سه مورد نخست، عملکردهای مشخص و مستقیمی دارند و غالباً بدون آرگومان‌های اضافی فراخوانی می‌شوند.

به‌عنوان نمونه، نحوه استفاده از دستور \cmd{reboot} جهت بازراه‌اندازی سیستم به شرح زیر است:

\begin{latin}
\begin{lstlisting}
[me@linuxbox ~]$ sudo reboot
\end{lstlisting}
\end{latin}

این فرمان به سیستم ابلاغ می‌کند که تمامی خدمات را به‌طور ایمن متوقف کرده و بلافاصله فرایند بوت مجدد را آغاز نماید.

در مقابل، دستور \cmd{shutdown} نسبت به دو دستور \cmd{halt} و \cmd{reboot} از انعطاف‌پذیری و کنترل بیشتری برخوردار است. این ابزار به شما اجازه می‌دهد تا نوع عملیات (توقف کامل یا بازراه‌اندازی) و همچنین بازه زمانی تأخیر برای اجرای دستور را دقیقاً تعیین کنید. این قابلیت به‌ویژه در محیط‌های چندکاربره (\en{Multi-user}) جهت اطلاع‌رسانی به کاربران فعال دربارهٔ وقوع یک رویداد قریب‌الوقوع و جلوگیری از از دست رفتن داده‌ها، بسیار حائز اهمیت است.

جهت خاموش کردن آنی سیستم، از دستور \cmd{shutdown} در سطح دسترسی «ابرکاربر» (\en{Superuser}) به شرح زیر استفاده می‌شود:

\begin{latin}
\begin{lstlisting}
[me@linuxbox ~]$ sudo shutdown -h now
\end{lstlisting}
\end{latin}

\begin{itemize}
  \item \cmd{-h}: این گزینه مخفف \en{halt} است و به سیستم فرمان می‌دهد که پس از خاتمه تمامی فرآیندها، سخت‌افزار را متوقف کند.
  \item \cmd{now}: پارامتر زمانی است که به سیستم ابلاغ می‌کند عملیات خاموشی را بدون درنگ آغاز نماید.
\end{itemize}

همچنین برای بازراه‌اندازی (\en{Reboot}) فوری سیستم، دستور زیر مورد استفاده قرار می‌گیرد:

\begin{latin}
\begin{lstlisting}
[me@linuxbox ~]$ sudo shutdown -r now
\end{lstlisting}
\end{latin}

\begin{itemize}
  \item \cmd{-r}: این گزینه مخفف \en{reboot} است و سیستم را پس از طی کردن مراحل توقف ایمن، مجدداً راه‌اندازی می‌کند.
\end{itemize}

شایان ذکر است که در سیستم‌های مدرن لینوکسی، استفاده از مقدار \cmd{now} معادل زمان‌بندی \cmd{+0} است.

دستور \cmd{shutdown} انعطاف‌پذیری بالایی در تعیین زمان‌بندی عملیات ارائه می‌دهد؛ به طوری که می‌توان تأخیر در اجرا را بر حسب دقیقه یا در یک ساعت مشخص از شبانه‌روز تنظیم کرد. بلافاصله پس از صدور این فرمان، یک پیام هشدار عمومی (\en{Broadcast}) برای تمامی کاربران فعال در سیستم ارسال می‌شود تا آن‌ها را از رویداد پیش‌رو مطلع ساخته و فرصت کافی برای ذخیره‌سازی داده‌ها فراهم کند. جهت بررسی جزئیات فنی و پارامترهای تکمیلی، می‌توانید به صفحه راهنمای سیستم (\cmd{man shutdown}) مراجعه کنید.

\section{ابزارهای تکمیلی پایش و مدیریت پردازش‌ها}

نظارت مستمر بر فرآیندها از کلیدی‌ترین وظایف مدیریت سیستم (\en{System Administration}) محسوب می‌شود و برای این منظور ابزارهای متنوعی توسعه یافته است. در جدول زیر، برخی از کاربردی‌ترین دستورات پایش فرآیندها معرفی شده‌اند:

\begin{table}[htbp]
\centering
\caption{دستورات تکمیلی مدیریت پردازش‌ها}
\label{tab:process-tools}
\begin{tabularx}{\textwidth}{l X}
\toprule
\textbf{دستور} & \textbf{توضیحات فنی} \\
\midrule
\cmd{pstree} & فرآیندهای در حال اجرا را در قالب یک ساختار درختی نمایش می‌دهد و روابط والد-فرزند (\en{Parent-Child}) میان آن‌ها را ترسیم می‌کند. این نمای بصری، درک سلسله‌مراتب اجرایی سیستم را تسهیل می‌سازد. \\
\addlinespace
\cmd{vmstat} & گزارشی جامع از وضعیت منابع سیستم شامل حافظهٔ فیزیکی، فضای مبادله (\en{Swap}) و نرخ ورودی/خروجی (\en{I/O}) ارائه می‌دهد. برای به‌روزرسانی لحظه‌ای خروجی، می‌توان یک بازهٔ زمانی (بر حسب ثانیه) تعیین کرد (برای نمونه، \cmd{vmstat 5}). جهت خروج از این پایش، از میان‌بر \cmd{Ctrl+c} استفاده می‌شود. \\
\addlinespace
\cmd{xload} & ابزاری گرافیکی که نمودار بار سیستم (\en{System Load}) را به‌صورت لحظه‌ای در یک پنجرهٔ مستقل نمایش می‌دهد و تغییرات فشار روی پردازنده را در بازه‌های زمانی مختلف ترسیم می‌کند. \\
\addlinespace
\cmd{tload} & عملکردی مشابه \cmd{xload} دارد، با این تفاوت که نمودار بار سیستم را مستقیماً درون محیط ترمینال و با استفاده از کاراکترهای متنی ترسیم می‌کند. توقف این نمایش نیز با کلیدهای \cmd{Ctrl+c} امکان‌پذیر است. \\
\bottomrule
\end{tabularx}
\end{table}

\section{جمع‌بندی}

سیستم‌های عامل مدرن و در رأس آن‌ها لینوکس، جهت مدیریت هم‌زمان وظایف متعدد، از مکانیزم‌های پیشرفته‌ای برای پایش و کنترل فرآیندها بهره می‌برند. لینوکس به عنوان پیشروترین سیستم‌عامل در حوزه سرور، مجموعه‌ای کامل از ابزارهای قدرتمند خط فرمان (\en{CLI}) را برای مدیریت دقیق چرخه حیات پردازش‌ها در اختیار مدیران سیستم قرار داده است.

هرچند رابط‌های گرافیکی (\en{GUI}) متعددی برای پایش سیستم توسعه یافته‌اند، اما ابزارهای خط فرمان به دلیل سرعت پاسخ‌گویی بالا و کمینه‌سازی مصرف منابع، همچنان اولویت اول متخصصان هستند. از دیدگاه مهندسی، یک ابزار نظارتی نباید خود به عاملی برای تحمیل بار اضافی بر سیستم تبدیل شود. از این رو، مدیران با‌تجربه ابزارهای متنی را به جایگزین‌های گرافیکی ترجیح می‌دهند؛ چرا که این ابزارها در تشخیص ریشه‌ای اختلالات و رفع گلوگاه‌های عملکردی، کارایی و دقت فنی بالاتری ارائه می‌دهند.